\documentclass[conference]{IEEEtran}

\usepackage{cite}
\usepackage{amsmath,amssymb}
\usepackage{graphicx}
\usepackage{booktabs}
\usepackage{multirow}
\usepackage{url}

\title{Beyond Accuracy: Leakage-Aware and Cross-Dataset AI-Generated Text Detection Using Lexical and Stylometric Features}

\author{
\IEEEauthorblockN{Salma Hesham Salem}
\IEEEauthorblockA{
Department of Computer Science\\
}
}

\begin{document}
\maketitle

\begin{abstract}
Detecting whether content was written by a human or generated by a large language model is often reported as a high-accuracy binary classification problem. However, high in-domain accuracy can be misleading when datasets contain label-revealing artifacts or when detectors are evaluated only on random splits from a single benchmark. This paper presents a leakage-aware and cross-dataset study of AI-generated content detection. We first audit the original Global AI vs Human Content dataset used in the pilot stage and show that it contains explicit AI markers, label-correlated metadata, and duplicate/conflicting examples. We therefore use it only as a leakage case study. The final experimental system compares four text-only detectors over SemEval-2024 Task 8 Subtask A, a balanced RAID subset, and an Arabic proof-of-concept dataset. The models include word TF-IDF with logistic regression, character TF-IDF with linear SVM, stylometric random forest, and a hybrid lexical-stylometric SVM. Within-dataset results are strong: the hybrid model reaches 0.9443 F1 on SemEval, while character TF-IDF reaches 0.9528 F1 on RAID. However, cross-dataset transfer is substantially weaker, with the best SemEval-to-RAID F1 at 0.7887 and the best RAID-to-SemEval F1 at 0.5977. These results show that robust AI-authorship detection requires leakage audits and cross-benchmark evaluation, not only high random-split scores.
\end{abstract}

\begin{IEEEkeywords}
AI-generated text detection, stylometry, leakage audit, cross-dataset evaluation, SemEval, RAID, Arabic NLP
\end{IEEEkeywords}

\section{Introduction}
Large language models such as GPT, Claude, Gemini, and open-weight instruction-tuned models have made machine-generated text increasingly fluent. This creates a practical need for detectors that can distinguish human-written content from generated content in educational, journalistic, and platform-moderation settings. A common approach is to treat the task as binary classification, where human-written text is assigned label 0 and AI-generated text is assigned label 1.

Despite strong benchmark results, AI-generated text detection remains fragile. Dataset artifacts can make the task artificially easy: generated rows may contain explicit markers, metadata may reveal labels, and templated examples may appear in both train and test splits. Even when these issues are removed, a detector trained on one benchmark may fail on another because domains, generators, decoding methods, prompts, and human writing styles differ.

This work is framed around that problem. The project began with an original Global AI vs Human Content dataset. Instead of treating it as the final benchmark, we audit it and use it as a motivation for a stricter experimental design. The final study uses SemEval-2024 Task 8 Subtask A as the main English benchmark, RAID as a robustness benchmark, and a KFUPM-JRCAI Arabic generated-text dataset as a multilingual proof-of-concept.

The contributions are:
\begin{itemize}
    \item a leakage-aware audit of the original dataset, showing why high pilot scores are not sufficient;
    \item a four-model comparison over lexical, character-level, stylometric, and hybrid features;
    \item within-dataset and cross-dataset evaluation across SemEval and RAID;
    \item an Arabic proof-of-concept using the same text-only detection pipeline.
\end{itemize}

\section{Related Work}
AI-generated text detection has been approached using neural classifiers, statistical features, perplexity-based methods, and stylometric analysis. Early authorship attribution work showed that character n-grams and writing-style statistics are strong signals for identifying authors and genres \cite{stamatatos2009survey}. Modern AI-text detection extends this idea by searching for machine signatures in lexical diversity, repetition, punctuation patterns, sentence length, and token distribution.

SemEval-2024 Task 8 introduced shared-task benchmarks for detecting machine-generated text across settings and languages \cite{wang2024semeval}. RAID was proposed as a robustness benchmark containing generations from multiple models, domains, decoding strategies, and adversarial attacks \cite{dugan2024raid}. These benchmarks are important because they move evaluation beyond a single small dataset.

Recent work also emphasizes that AI-text detectors often fail under distribution shift \cite{sadasivan2023can}. A model may perform well on the generator or prompt style seen during training but degrade on unseen generators, new domains, or adversarially modified outputs. This motivates the cross-dataset protocol used in this paper.

\section{Initial Dataset Audit}
The project originally used a local Global AI vs Human Content dataset with 20,000 rows. It includes columns such as content, label, type, source, prompt, topic, word count, character count, AI model, language, complexity score, and code flags. The audit found three major risks.

First, many AI-generated rows contain explicit textual markers such as ``AI-generated''. Second, metadata columns such as \texttt{source} and \texttt{ai\_model} are label-correlated and can reveal the target if used directly or indirectly. Third, normalized duplicates and conflicting examples exist, meaning a random split may overestimate generalization.

For these reasons, this dataset is used as Dataset 0: a pilot and leakage-audit case study. It is not treated as the main final benchmark. Table~\ref{tab:audit} summarizes its role.

\begin{table}[t]
\centering
\caption{Dataset 0 Audit Summary}
\label{tab:audit}
\begin{tabular}{ll}
\toprule
Item & Observation \\
\midrule
Original size & 20,000 rows \\
Prepared size & 8,242 deduplicated rows \\
Problems & AI markers, metadata leakage, duplicates \\
Within-dataset F1 & 0.9994--1.0000 \\
Final role & Leakage audit / pilot only \\
\bottomrule
\end{tabular}
\end{table}

The old dataset was also trained and evaluated for comparison. Its within-dataset scores are nearly perfect, but old-to-new transfer is poor. This supports the decision to avoid using it as the primary benchmark.

\section{Datasets}
\subsection{SemEval English}
SemEval-2024 Task 8 Subtask A is used as the main English binary benchmark. The prepared English monolingual dataset contains 124,552 cleaned examples after missing-text removal, whitespace normalization, and duplicate removal. The class distribution is 65,646 human and 58,906 AI examples.

\subsection{RAID English}
RAID is used as a second English benchmark for robustness. Because the full dataset is large, a controlled streamed train subset is prepared and balanced. The final cleaned RAID subset contains 3,532 examples: 1,766 human and 1,766 AI examples. Metadata such as domain, generator, attack, and decoding is retained only for analysis.

\subsection{Arabic Proof-of-Concept}
The Arabic proof-of-concept uses KFUPM-JRCAI Arabic generated abstracts. The wide-format dataset is transformed into long binary format with original abstracts labeled human and model-generated abstracts labeled AI. The final Arabic set contains 36,523 examples. It is highly imbalanced, so it is reported as a proof-of-concept rather than a final multilingual benchmark.

\section{Methodology}
All experiments use text-only input. Metadata fields such as generator, domain, attack, source, prompt, dataset name, and AI model name are excluded from model features.

Four models are compared:
\begin{itemize}
    \item \textbf{M1 Word TF-IDF + Logistic Regression}: word unigrams/bigrams with balanced logistic regression.
    \item \textbf{M2 Character TF-IDF + Linear SVM}: character 3--5 grams with a linear SVM.
    \item \textbf{M3 Stylometric Random Forest}: handcrafted features such as length, lexical diversity, punctuation rates, sentence statistics, digit ratio, uppercase ratio, and repetition.
    \item \textbf{M4 Hybrid TF-IDF + Stylometric SVM}: word TF-IDF, character TF-IDF, and stylometric features concatenated and classified with a linear SVM.
\end{itemize}

English cleaning is intentionally light: missing and empty text are removed, duplicates are removed, and whitespace is normalized. Punctuation and casing are preserved because they are meaningful stylometric signals. Arabic text is normalized by removing diacritics and tatweel, normalizing Alef variants, normalizing \textit{ى} to \textit{ي}, and collapsing whitespace.

\section{Experimental Setup}
Each dataset uses a single stratified train/test split with random seed 42 and test size 0.2. All models use the same split within each dataset. The main evaluation metrics are accuracy, precision, recall, F1, and ROC-AUC.

Three evaluation views are reported:
\begin{enumerate}
    \item within-dataset SemEval and RAID evaluation;
    \item cross-dataset transfer between SemEval and RAID;
    \item old-dataset comparison to test whether the pilot dataset generalizes.
\end{enumerate}

\section{Results}
\subsection{Within-Dataset English Results}
Table~\ref{tab:english} reports the main English results. On SemEval, the hybrid model performs best with F1 = 0.9443. On RAID, the character-level SVM performs best with F1 = 0.9528, slightly ahead of the hybrid model.

\begin{table}[t]
\centering
\caption{Within-Dataset English Results}
\label{tab:english}
\begin{tabular}{llcc}
\toprule
Dataset & Best Model & Accuracy & F1 \\
\midrule
SemEval & M4 Hybrid & 0.9475 & 0.9443 \\
RAID & M2 Char TF-IDF SVM & 0.9533 & 0.9528 \\
\bottomrule
\end{tabular}
\end{table}

The result should not be interpreted as ``hybrid is always best.'' Instead, it shows that different benchmarks reward different signals. SemEval benefits from combining lexical, character, and stylometric features, while RAID is very well captured by character n-grams.

\subsection{Old Dataset Results}
The old dataset gives near-perfect within-dataset performance. M1, M2, and M3 reach 1.0000 F1, while M4 reaches 0.9994 F1. These numbers are not used as final evidence of detector quality because the audit identified leakage risks. They instead support the paper's central warning: strong random-split results can be misleading.

\section{Cross-Dataset Generalization}
Cross-dataset evaluation is the strongest finding. Table~\ref{tab:cross} shows that transfer is weaker than within-dataset testing. SemEval-to-RAID transfer reaches a best F1 of 0.7887 using M2. RAID-to-SemEval transfer is substantially worse, with a best F1 of 0.5977 using M1.

\begin{table}[t]
\centering
\caption{Best Cross-Dataset Results}
\label{tab:cross}
\begin{tabular}{llcc}
\toprule
Train $\rightarrow$ Test & Best Model & Accuracy & F1 \\
\midrule
SemEval $\rightarrow$ RAID & M2 Char TF-IDF SVM & 0.8091 & 0.7887 \\
RAID $\rightarrow$ SemEval & M1 Word TF-IDF LR & 0.5721 & 0.5977 \\
\bottomrule
\end{tabular}
\end{table}

The old dataset comparison reinforces this point. Models trained on the old dataset generalize poorly to SemEval and RAID. For example, old-to-SemEval F1 is as low as 0.0364 for M2 and only 0.6154 for the best old-trained model. Old-to-RAID F1 is also weak, with the best result at 0.5042. This confirms that the old dataset is useful as an audit case, not as a robust benchmark.

\section{Arabic Proof-of-Concept}
The Arabic proof-of-concept results are strong: M2 reaches 0.9937 F1 and M4 reaches 0.9940 F1. However, the Arabic dataset is imbalanced because each original text is paired with several generated variants. Therefore, these results are reported as evidence that the pipeline can be adapted to Arabic, not as a complete Arabic benchmark. Future work should add balanced Arabic datasets, more domains, and cross-dataset Arabic evaluation.

\section{Limitations}
This study has several limitations. First, RAID is represented by a streamed balanced subset rather than the full benchmark. Second, the Arabic dataset is imbalanced and domain-specific. Third, the classical models do not include transformer fine-tuning, which may improve within-dataset performance but can also overfit to benchmark artifacts. Fourth, the old dataset audit identifies leakage risks but does not attempt to fully repair the dataset into a publication-grade benchmark.

\section{Conclusion}
This paper argues that AI-generated text detection should be evaluated beyond random-split accuracy. The original dataset produces near-perfect scores but shows leakage risks and poor transfer, so it is used as a leakage-audit case study. On stronger English benchmarks, the hybrid model performs best on SemEval and character TF-IDF performs best on RAID. Cross-dataset evaluation reveals a substantial robustness gap, especially from RAID to SemEval. The Arabic proof-of-concept shows that the pipeline can be extended to non-English text, but more balanced multilingual benchmarks are needed. Overall, the main contribution is a leakage-aware, cross-dataset evaluation pipeline for AI-authorship detection.

\bibliographystyle{IEEEtran}
\bibliography{references}

\end{document}
